\section{Related work}
\label{sec:related}

\paragraph{Analytic choices and the multiverse.} A multiverse analysis computes a result under every reasonable
combination of processing choices and reports the whole set \cite{steegen2016multiverse}; a specification curve orders
the results and shows which choices drive them \cite{simonsohn2020specification}. In educational measurement, a
specification curve over scaling decisions for PISA showed that model uncertainty moved country means about as much as
the sampling of students \cite{robitzsch2022multiverse}; we take from it the idea of setting definitional and sampling
variation side by side. In educational data mining, time-on-task estimation changed regression results and their
significance \cite{kovanovic2015time}, two operationalisations of wheel spinning agreed on less than half of the cases
\cite{zhang2019wheel}, and predictive models of student success were examined as a multiverse
\cite{tang2025multiverse}. For behavioural data closer to ours, a multiverse of cleaning and analysis procedures for
eye movements in reading found the word-frequency effect almost always present but varying substantially in size; it
did not vary the algorithm that parses samples into fixations \cite{godwin2025multiverse}. Pupillometry has been
analysed in the same way \cite{calignano2023pupillometry}, and eye-movement detectors have been compared in the timing
of the events they label \cite{nir2026systematic}. Adjusting for a covariate can change the sign of an association, the
pattern known as suppression \cite{mackinnon2000equivalence}. Our analysis varies the cutting of events itself and
reads, beyond the size of an effect, its direction, the comparison between datasets and the structure within writers.

\paragraph{Arbitrary pause thresholds and the answers to them.} In CS1 keystroke data, inactivity thresholds decide time
on task \cite{hart2023time,hart2025time,edwards2022engagement}, the frequency of pauses in several length bands relates to
performance \cite{shrestha2022pausing}, and a multi-cohort replication found that pausing at chunk boundaries remained
negatively related to exam scores after controlling for typing fluency and general hesitation, with a strength that
changed when the pauses were split at 2~s \cite{koutcheme2026keystroke}. In writing, fixed thresholds such as 2~s are
criticised as having no theoretical basis \cite{torrance2024methods,roeser2024disfluencies}, and alternatives have been
proposed: thresholds personalised to the writer \cite{chukharev2014pauses}, measures grounded in the transition
\cite{hall2024constructing}, pauses classified by linguistic location \cite{medimorec2017pauses}, mixture models of
pause durations \cite{baaijen2012keystroke,almond2012preliminary,li2021mixture,litvinova2024investigating}, personalised
thresholds proportional to the writer's within-word speed \cite{guo2020scenario}, state models
that avoid thresholds \cite{garcesmanzanera2026modeling}, and control for typing skill \cite{alves2007typing}. Each
answer still makes a choice---of a threshold, a model or a control---so whether its choice changes the conclusions can
be checked with the same kind of diagnostic. We add that check rather than another answer.

\paragraph{Pauses and essay scores.} In single studies, measures of pause duration and latency correlate negatively
with essay scores, from about $-0.15$ to $-0.43$ \cite{li2026bayesian,zhang2015process,schaller2026keyscore,li2021mixture};
text length correlates with the score at 0.73 in a large corpus, and the number of pauses can correlate
positively with the score through length \cite{schaller2026keyscore}. The direction differs with the location of the
pause \cite{zhang2015process}, the stage of writing \cite{xu2014pauses} and the task condition
\cite{schaller2026keyscore}, and differing thresholds are said to make results hard to compare \cite{xu2014pauses}.
Within our search (S7) we found no study that varied the value of the threshold and reported whether the direction of
the association with the score changed.

\paragraph{Levels, stability and dispersion of keystroke measures.} Keystroke process features vary across testing
occasions \cite{deane2015feasibility} and across a student's portfolio \cite{babalola2026stability}; the generalisability
of pause times in sentence production for distinguishing writers has been estimated
\cite{meulemans2022generalizability}; mobile typing has been modelled with time-of-day and mood effects
\cite{vesel2020biaffect}. Counts of pauses are often modelled as Poisson with a check for overdispersion
\cite{gullberg2025deceit}. Within the scope of our search (S7), we found no study that decomposed stall length into
three or more levels or reported the dispersion of stall counts within writers across definitions.
